\documentclass[UTF8,fontset=none,11pt,a4paper]{ctexart}
\usepackage[left=21mm,right=21mm,top=27mm,bottom=22mm,headheight=26pt,headsep=6mm]{geometry}
\usepackage{fontspec}
\setmainfont{texgyretermes-regular.otf}[BoldFont=texgyretermes-bold.otf,ItalicFont=texgyretermes-italic.otf]
\setCJKmainfont{Songti.ttc}[Path=/System/Library/Fonts/Supplemental/,FontIndex=6,AutoFakeBold=2]
\setCJKsansfont{Songti.ttc}[Path=/System/Library/Fonts/Supplemental/,FontIndex=6,AutoFakeBold=2]
\usepackage[table]{xcolor}
\definecolor{TableHead}{HTML}{EEF0ED}
\definecolor{Muted}{HTML}{666C68}
\usepackage{booktabs,longtable,array,fancyhdr,needspace,xurl}
\usepackage[unicode,hidelinks]{hyperref}
\urlstyle{same}
\linespread{1.32}
\setlength{\parindent}{2em}
\setlength{\parskip}{3pt}
\setlength{\LTpre}{8pt}\setlength{\LTpost}{10pt}
\ctexset{section={format=\fontsize{15}{20}\selectfont\bfseries,beforeskip=18pt,afterskip=8pt},subsection={format=\fontsize{12}{16}\selectfont\bfseries,beforeskip=12pt,afterskip=6pt},subsubsection={format=\normalsize\bfseries,beforeskip=9pt,afterskip=4pt}}
\pagestyle{fancy}\fancyhf{}
\renewcommand{\headrulewidth}{.3pt}
\renewcommand{\footrulewidth}{0pt}
\fancyhead[L]{\parbox[b]{\linewidth}{\fontsize{8}{10}\selectfont\color{Muted}退休账户转移：五年持有与两层费用·基金F招募说明书}}
\fancyfoot[C]{\small\color{Muted}\thepage}
\hypersetup{pdftitle={基金F招募说明书},pdfsubject={仅供模拟使用·退休账户转移：五年持有与两层费用·03-02},pdfauthor={金融决策模拟材料}}
\begin{document}
{\noindent\fontsize{20}{25}\selectfont\bfseries 基金F招募说明书\par}
\vspace{7pt}{\noindent\small\color{Muted}仅供模拟使用·退休账户转移：五年持有与两层费用·03-02\par}\vspace{10pt}
P类及Q类份额


2026年3月1日


\section*{投资目标}
基金投资目标为保存资本、取得当前收入以及资本和收入的长期增长。基金董事会目前无意改变目标，但可在提前60日书面通知持有人后不经持有人批准变更。


\section*{基金费用及支出}
下表列示购买、持有及出售本基金份额可能承担的费用及支出。您还可能承担支付给金融中介的经纪佣金及其他费用，这些费用不在下表及费用示例中列示。P类和Q类份额投资于同一证券组合，各类别分别承担其应适用的费用。


\Needspace{5\baselineskip}
\begingroup\small\setlength{\tabcolsep}{4pt}\renewcommand{\arraystretch}{1.35}
\begin{longtable}{>{\centering\arraybackslash}m{0.45163\dimexpr\linewidth-6\tabcolsep\relax}>{\centering\arraybackslash}m{0.27418\dimexpr\linewidth-6\tabcolsep\relax}>{\centering\arraybackslash}m{0.27418\dimexpr\linewidth-6\tabcolsep\relax}}
\toprule
\rowcolor{TableHead} \centering\arraybackslash \textbf{持有人直接承担\mbox{的费用}} & \centering\arraybackslash \textbf{P类别} & \centering\arraybackslash \textbf{Q类别} \\
\midrule\endfirsthead
\toprule
\rowcolor{TableHead} \centering\arraybackslash \textbf{持有人直接承担\mbox{的费用}} & \centering\arraybackslash \textbf{P类别} & \centering\arraybackslash \textbf{Q类别} \\
\midrule\endhead
\midrule\endfoot
\bottomrule\endlastfoot
\centering\arraybackslash 最大购买\mbox{销售费} & \centering\arraybackslash 无 & \centering\arraybackslash 无 \\ \addlinespace[3pt]
\centering\arraybackslash 最大递延\mbox{销售费} & \centering\arraybackslash 无 & \centering\arraybackslash 无 \\ \addlinespace[3pt]
\centering\arraybackslash 赎回或\mbox{转换费} & \centering\arraybackslash 无 & \centering\arraybackslash 无 \\
\end{longtable}\endgroup
\Needspace{5\baselineskip}
\begingroup\small\setlength{\tabcolsep}{4pt}\renewcommand{\arraystretch}{1.35}
\begin{longtable}{>{\centering\arraybackslash}m{0.52444\dimexpr\linewidth-6\tabcolsep\relax}>{\centering\arraybackslash}m{0.23778\dimexpr\linewidth-6\tabcolsep\relax}>{\centering\arraybackslash}m{0.23778\dimexpr\linewidth-6\tabcolsep\relax}}
\toprule
\rowcolor{TableHead} \centering\arraybackslash \textbf{年度基金运营费用，以投资净资产百分\mbox{比表示}} & \centering\arraybackslash \textbf{P类别} & \centering\arraybackslash \textbf{Q类别} \\
\midrule\endfirsthead
\toprule
\rowcolor{TableHead} \centering\arraybackslash \textbf{年度基金运营费用，以投资净资产百分\mbox{比表示}} & \centering\arraybackslash \textbf{P类别} & \centering\arraybackslash \textbf{Q类别} \\
\midrule\endhead
\midrule\endfoot
\bottomrule\endlastfoot
\centering\arraybackslash \mbox{管理费} & \centering\arraybackslash 0.21\% & \centering\arraybackslash 0.21\% \\ \addlinespace[3pt]
\centering\arraybackslash 分销或服务费（12b-1） & \centering\arraybackslash 0.75\% & \centering\arraybackslash 无 \\ \addlinespace[3pt]
\centering\arraybackslash 其\mbox{他费用} & \centering\arraybackslash 0.37\% & \centering\arraybackslash 0.04\% \\ \addlinespace[3pt]
\centering\arraybackslash 年度基金运营费\mbox{用合计} & \centering\arraybackslash 1.33\% & \centering\arraybackslash 0.25\% \\
\end{longtable}\endgroup
费用表仅列基金自身费用。通过收费投资渠道持有时，渠道另收的持续服务费用不由基金支付。


\subsection*{费用示例}
本示例旨在帮助您比较投资本基金与投资其他共同基金的成本。


假设您向本基金投资10,000美元，在下表所列期间结束时赎回或继续持有全部份额；另假设您的投资每年取得5\%的回报，基金运营费用保持不变。买卖Q类份额时可能须另付经纪佣金，该佣金不包括在本示例内。虽然您的实际成本可能高于或低于下列金额，在这些假设下，您承担的成本如下：


\Needspace{5\baselineskip}
\begingroup\small\setlength{\tabcolsep}{4pt}\renewcommand{\arraystretch}{1.35}
\begin{longtable}{>{\centering\arraybackslash}m{0.31266\dimexpr\linewidth-6\tabcolsep\relax}>{\centering\arraybackslash}m{0.37468\dimexpr\linewidth-6\tabcolsep\relax}>{\centering\arraybackslash}m{0.31266\dimexpr\linewidth-6\tabcolsep\relax}}
\toprule
\rowcolor{TableHead} \centering\arraybackslash \textbf{持\mbox{有期间}} & \centering\arraybackslash \textbf{P类} & \centering\arraybackslash \textbf{Q类} \\
\midrule\endfirsthead
\toprule
\rowcolor{TableHead} \centering\arraybackslash \textbf{持\mbox{有期间}} & \centering\arraybackslash \textbf{P类} & \centering\arraybackslash \textbf{Q类} \\
\midrule\endhead
\midrule\endfoot
\bottomrule\endlastfoot
\centering\arraybackslash 1年 & \centering\arraybackslash 135美元 & \centering\arraybackslash 26美元 \\ \addlinespace[3pt]
\centering\arraybackslash 3年 & \centering\arraybackslash 421美元 & \centering\arraybackslash 80美元 \\ \addlinespace[3pt]
\centering\arraybackslash 5年 & \centering\arraybackslash 729美元 & \centering\arraybackslash 141美元 \\ \addlinespace[3pt]
\centering\arraybackslash 10年 & \centering\arraybackslash 1,601美元 & \centering\arraybackslash 318美元 \\
\end{longtable}\endgroup
组合周转。基金买卖证券时支付佣金等交易成本。较高周转可能表示较高交易成本，并在应税账户产生较高税负。这些成本不在年度基金运营费用中，但影响基金投资结果。最近财务年度组合周转率为平均组合价值的106\%。


\section*{主要投资策略}
基金以平衡方法投资广泛证券，包括普通股和投资级债券。投资级为指定评级机构Baa3、BBB-或更高等级，或未评级但投资顾问判断具有相当质量的证券。基金也投资美国政府及联邦机构发行和担保的证券，并可将部分资产投资于美国以外的普通股、债券及其他证券，大部分股票有支付股息的历史。


正常情况下，至少50\%的资产价值投资普通股，至少25\%投资包括货币市场证券在内的债务证券。重点为中型及较大型公司，但不限于某一市值规模。债务证券不受特定到期日或久期限制。久期用于衡量价格对利率变化的敏感性，久期越长，通常越敏感。


基金采用多组合经理制度，将资产分为不同部分，由各经理独立管理。投资顾问依专业判断投资其认为估值有吸引力、具有良好长期投资机会的证券；不再相对有吸引力时可以出售。


\section*{主要风险}
投资基金可能亏损，投资期限较短时损失可能性可能更大。


\textbf{市场条件。}基金持有的股票、债券及其他证券价格和收入可因整体经济、行业或公司情况、市场变化、政治社会经济不稳定、政府和央行应对、公共债务与赤字、通胀，以及汇率、利率和商品价格变化而下跌，有时迅速且难以预料。


世界经济及金融市场紧密联系。贸易安排、战争、恐怖主义、网络安全事件、自然灾害、公共卫生紧急情况、银行失败及政府应对可影响全球。即使基金不投资受影响国家的发行人，其他地区变化仍可影响价值和流动性。


\textbf{发行人风险。}需求下降、管理表现不佳、重大诉讼、调查或争议、财务状况或评级变化、监管及竞争环境，以及并购或资产处置等战略行动，可能降低证券价格和收入。个别证券可受行业影响，单一发行人也可能影响行业或整个市场。


\textbf{增长型股票。}普通股及其他股票型证券，包括优先股、可转换优先股和可转换债券，相比其他投资可有更大波动及损失可能。收入型股票。公司分红政策及可用于分红的资本资源变化，可能减少证券价值及基金收入。


\textbf{债务工具。}价格和收入受利率、到期日和信用质量影响。利率上升通常导致债券下跌，亦可能降低提前偿还速度并延长存续期；大规模卖出可降低价格和流动性并增加基金赎回。利率下降可导致提前赎回、偿还或再融资，基金须以较低收益再投资。较长期限通常更敏感。


\textbf{信用风险。}发行人或担保人信用实力可能减弱或被视为减弱，也可能未及时付本息并违约。降级或违约可降低基金份额价值。较低信用质量债券通常利率较高，价格波动更大。评级是部分判断依据，但不保证信用质量。


\textbf{美国政府证券。}政府直接发行或完全担保的证券，仅在持有至到期时保证按期支付本息，其市场价值仍随利率和政府评级变化。即使有政府完全信用支持，特殊情况也可能妨碍或延迟付款。非由政府完全信用支持的政府支持实体或联邦机构证券，不属于美国政府直接发行或担保。


\textbf{抵押及资产支持证券。}基础资产可能是住房贷款、消费者贷款或应收款。除一般债务风险外，借款人付款速度变化影响价值。利率下降时提前还款增多，须以较低收益再投资；利率上升时偿还较慢，资金不能及时再投资高收益工具。借款人可能违约，担保房地产可能不足以在止赎后偿债。其他资产支持证券面临类似风险。


\textbf{流动性风险。}部分持仓在市场动荡时难以或无法出售。没有活跃市场、法律或合同转售限制，或做市人减少，都可能降低流动性。价格可能波动或难以确定。出售可能延误且成本增加，基金可能无法在需要时出售，或只能亏损出售。


美国以外投资。政治社会经济和市场变化、地区冲突及战争、汇率变化、国有化、外汇限制、制裁、税收和价格控制，可降低价值。外国市场可能波动更大、流动性更低，适用不同监管、法律、会计、审计及披露要求，估值更困难。外国税收和交易结算延误也可影响结果，新兴市场风险可能更高。


\section*{基金投资业绩}
下表列示基金不同期间的平均年度总回报，并与证券市场的广泛衡量指标及反映基金投资范围的其他市场指标比较。这些资料有助于说明投资本基金的风险。基金过往投资业绩，无论税前还是税后，均不能用于预测未来投资业绩。


\subsection*{平均年度总回报}
截至2025年12月31日的期间。基金业绩包括股息及资本利得分配的再投资；除另有说明外，反映列示期间内有效的费用减免及费用补偿。P类和Q类下表回报均为税前回报。


\Needspace{5\baselineskip}
\begingroup\small\setlength{\tabcolsep}{4pt}\renewcommand{\arraystretch}{1.35}
\begin{longtable}{>{\centering\arraybackslash}m{0.1524\dimexpr\linewidth-12\tabcolsep\relax}>{\centering\arraybackslash}m{0.23799\dimexpr\linewidth-12\tabcolsep\relax}>{\centering\arraybackslash}m{0.1524\dimexpr\linewidth-12\tabcolsep\relax}>{\centering\arraybackslash}m{0.1524\dimexpr\linewidth-12\tabcolsep\relax}>{\centering\arraybackslash}m{0.1524\dimexpr\linewidth-12\tabcolsep\relax}>{\centering\arraybackslash}m{0.1524\dimexpr\linewidth-12\tabcolsep\relax}}
\toprule
\rowcolor{TableHead} \centering\arraybackslash \textbf{份\mbox{额类别}} & \centering\arraybackslash \textbf{开始运\mbox{作日期}} & \centering\arraybackslash \textbf{1年} & \centering\arraybackslash \textbf{5年} & \centering\arraybackslash \textbf{10年} & \centering\arraybackslash \textbf{成\mbox{立以来}} \\
\midrule\endfirsthead
\toprule
\rowcolor{TableHead} \centering\arraybackslash \textbf{份\mbox{额类别}} & \centering\arraybackslash \textbf{开始运\mbox{作日期}} & \centering\arraybackslash \textbf{1年} & \centering\arraybackslash \textbf{5年} & \centering\arraybackslash \textbf{10年} & \centering\arraybackslash \textbf{成\mbox{立以来}} \\
\midrule\endhead
\midrule\endfoot
\bottomrule\endlastfoot
\centering\arraybackslash P类 & \centering\arraybackslash 2002年5月21日 & \centering\arraybackslash 17.58\% & \centering\arraybackslash 8.73\% & \centering\arraybackslash 8.95\% & \centering\arraybackslash 7.22\% \\ \addlinespace[3pt]
\centering\arraybackslash Q类 & \centering\arraybackslash 2017年1月27日 & \centering\arraybackslash 18.84\% & \centering\arraybackslash 9.92\% & \centering\arraybackslash \mbox{不适用} & \centering\arraybackslash 10.08\% \\
\end{longtable}\endgroup
\Needspace{5\baselineskip}
\begingroup\small\setlength{\tabcolsep}{4pt}\renewcommand{\arraystretch}{1.35}
\begin{longtable}{>{\centering\arraybackslash}m{0.51504\dimexpr\linewidth-8\tabcolsep\relax}>{\centering\arraybackslash}m{0.16165\dimexpr\linewidth-8\tabcolsep\relax}>{\centering\arraybackslash}m{0.16165\dimexpr\linewidth-8\tabcolsep\relax}>{\centering\arraybackslash}m{0.16165\dimexpr\linewidth-8\tabcolsep\relax}}
\toprule
\rowcolor{TableHead} \centering\arraybackslash \textbf{比\mbox{较指数}} & \centering\arraybackslash \textbf{1年} & \centering\arraybackslash \textbf{5年} & \centering\arraybackslash \textbf{10年} \\
\midrule\endfirsthead
\toprule
\rowcolor{TableHead} \centering\arraybackslash \textbf{比\mbox{较指数}} & \centering\arraybackslash \textbf{1年} & \centering\arraybackslash \textbf{5年} & \centering\arraybackslash \textbf{10年} \\
\midrule\endhead
\midrule\endfoot
\bottomrule\endlastfoot
\centering\arraybackslash 标普500指数 & \centering\arraybackslash 17.88\% & \centering\arraybackslash 14.42\% & \centering\arraybackslash 14.82\% \\ \addlinespace[3pt]
\centering\arraybackslash 60\%标普500指数／40\%彭博美国综\mbox{合指数} & \centering\arraybackslash 13.70\% & \centering\arraybackslash 8.47\% & \centering\arraybackslash 9.78\% \\ \addlinespace[3pt]
\centering\arraybackslash 彭博美国综\mbox{合指数} & \centering\arraybackslash 7.30\% & \centering\arraybackslash −0.36\% & \centering\arraybackslash 2.01\% \\
\end{longtable}\endgroup
指数未经主动管理，其回报包括股息或分配的再投资，未扣除销售费用、佣金、账户费用、运营支出或美国联邦所得税。


\subsection*{管理}
\textbf{投资顾问：}管理公司M。


本基金采用多组合经理制度，由不同组合经理分别管理基金资产的一部分。主要负责本基金组合管理的人员、管理本基金的起始年份及其在投资顾问机构的主要职务如下：


\Needspace{5\baselineskip}
\begingroup\small\setlength{\tabcolsep}{4pt}\renewcommand{\arraystretch}{1.35}
\begin{longtable}{>{\centering\arraybackslash}m{0.3936\dimexpr\linewidth-6\tabcolsep\relax}>{\centering\arraybackslash}m{0.18484\dimexpr\linewidth-6\tabcolsep\relax}>{\centering\arraybackslash}m{0.42156\dimexpr\linewidth-6\tabcolsep\relax}}
\toprule
\rowcolor{TableHead} \centering\arraybackslash \textbf{组\mbox{合经理}} & \centering\arraybackslash \textbf{管理本基金起\mbox{始年份}} & \centering\arraybackslash \textbf{主\mbox{要职务}} \\
\midrule\endfirsthead
\toprule
\rowcolor{TableHead} \centering\arraybackslash \textbf{组\mbox{合经理}} & \centering\arraybackslash \textbf{管理本基金起\mbox{始年份}} & \centering\arraybackslash \textbf{主\mbox{要职务}} \\
\midrule\endhead
\midrule\endfoot
\bottomrule\endlastfoot
\centering\arraybackslash 经理01 & \centering\arraybackslash 1999年 & \centering\arraybackslash 股票投资部门\mbox{合伙人} \\ \addlinespace[3pt]
\centering\arraybackslash 经理02 & \centering\arraybackslash 2018年 & \centering\arraybackslash 固定收益投资部门\mbox{合伙人} \\ \addlinespace[3pt]
\centering\arraybackslash 经理03，基金联\mbox{合总裁} & \centering\arraybackslash 2014年 & \centering\arraybackslash 股票投资部门\mbox{合伙人} \\ \addlinespace[3pt]
\centering\arraybackslash 经理04 & \centering\arraybackslash 2005年 & \centering\arraybackslash 股票投资部门\mbox{合伙人} \\ \addlinespace[3pt]
\centering\arraybackslash 经理05 & \centering\arraybackslash 2021年 & \centering\arraybackslash 股票投资部门\mbox{合伙人} \\ \addlinespace[3pt]
\centering\arraybackslash 经理06 & \centering\arraybackslash 2018年 & \centering\arraybackslash 股票投资部门\mbox{合伙人} \\ \addlinespace[3pt]
\centering\arraybackslash 经理07 & \centering\arraybackslash 2022年 & \centering\arraybackslash 固定收益投资部门\mbox{合伙人} \\ \addlinespace[3pt]
\centering\arraybackslash 经理08，基金联\mbox{合总裁} & \centering\arraybackslash 2016年 & \centering\arraybackslash 固定收益投资部门\mbox{合伙人} \\ \addlinespace[3pt]
\centering\arraybackslash 经理09 & \centering\arraybackslash 2020年 & \centering\arraybackslash 股票投资部门\mbox{合伙人} \\ \addlinespace[3pt]
\centering\arraybackslash 经理10 & \centering\arraybackslash 2020年 & \centering\arraybackslash 固定收益投资部门\mbox{合伙人} \\ \addlinespace[3pt]
\centering\arraybackslash 经理11 & \centering\arraybackslash 2016年 & \centering\arraybackslash 股票投资部门\mbox{合伙人} \\
\end{longtable}\endgroup
\subsection*{基金份额的购买及出售}
各类别开立账户的最低投资额通常为250美元，追加投资最低额为50美元。工资扣款退休计划账户、工资扣款储蓄计划账户及雇主设立的教育储蓄账户，初次或追加投资最低额为25美元。由基金过户代理人直接持有并提供服务的Q类账户，最低投资额为100万美元。


零售投资者可在营业日通过经销商或财务专业人员，或者以书面、电话、传真及网上方式出售或赎回份额。退休计划持有人应联系计划管理人或记录机构办理出售及赎回。


\subsection*{税务资料}
基金向您支付的股息和资本利得分配须缴纳联邦所得税，并可能须缴纳州及地方税；免税持有人或税收优惠账户除外。在税收优惠账户内持有份额时，日后从账户提取投资可能须纳税。


\subsection*{向经纪交易商及其他金融中介付款}
通过经纪交易商、银行或其他金融中介购买份额时，本基金、分销机构或其关联机构可能就基金份额销售及相关服务向该中介支付款项。此类付款可能影响中介及您的财务专业人员，使其推荐本基金而非其他投资，从而产生利益冲突。有关详情，可向财务专业人员查询或查阅中介机构的披露。


\section*{投资目标、策略及风险}
基金可持有现金及现金等价物，包括商业票据和美国政府、机构发行的短期证券。比例取决于市场、申购及赎回。在市场动荡等情况下可较大比例持有；临时防御时可不受限制投资上述工具。较高现金比例可降低上涨市场的投资结果，也可在下跌时减少损失并提供后续投资或赎回的流动性。


基金可投资同一顾问或关联方管理的中央基金，以有效取得货币市场、债券及其他资产类别的分散投资。中央基金只面向顾问及关联方管理的基金、投资工具和账户，不另收管理费，但基金仍承担按份额分配的其他费用。投入中央基金的资产表现取决于中央基金投资结果。


基金还可向经纪商、交易商及其他机构借出组合证券，并取得至少等于借出证券价值的现金或美国国债作担保。


基本分析可包括与公司管理层、员工、供应商、客户和竞争对手会面。投资顾问可考虑对发行人或证券价值具有实质影响的环境、社会及治理因素。考量取决于具体事实与情形，并非每项投资采取相同判断权重。


\section*{主要风险}
投资基金可能亏损，投资期限较短时损失可能性可能更大。


\textbf{市场条件。}基金持有的股票、债券及其他证券价格和收入可因整体经济、行业或公司情况、市场变化、政治社会经济不稳定、政府和央行应对、公共债务与赤字、通胀，以及汇率、利率和商品价格变化而下跌，有时迅速且难以预料。


世界经济及金融市场紧密联系。贸易安排、战争、恐怖主义、网络安全事件、自然灾害、公共卫生紧急情况、银行失败及政府应对可影响全球。即使基金不投资受影响国家的发行人，其他地区变化仍可影响价值和流动性。


\textbf{发行人风险。}需求下降、管理表现不佳、重大诉讼、调查或争议、财务状况或评级变化、监管及竞争环境，以及并购或资产处置等战略行动，可能降低证券价格和收入。个别证券可受行业影响，单一发行人也可能影响行业或整个市场。


\textbf{增长型股票。}普通股及其他股票型证券，包括优先股、可转换优先股和可转换债券，相比其他投资可有更大波动及损失可能。收入型股票。公司分红政策及可用于分红的资本资源变化，可能减少证券价值及基金收入。


\textbf{债务工具。}价格和收入受利率、到期日和信用质量影响。利率上升通常导致债券下跌，亦可能降低提前偿还速度并延长存续期；大规模卖出可降低价格和流动性并增加基金赎回。利率下降可导致提前赎回、偿还或再融资，基金须以较低收益再投资。较长期限通常更敏感。


\textbf{信用风险。}发行人或担保人信用实力可能减弱或被视为减弱，也可能未及时付本息并违约。降级或违约可降低基金份额价值。较低信用质量债券通常利率较高，价格波动更大。评级是部分判断依据，但不保证信用质量。


\textbf{美国政府证券。}政府直接发行或完全担保的证券，仅在持有至到期时保证按期支付本息，其市场价值仍随利率和政府评级变化。即使有政府完全信用支持，特殊情况也可能妨碍或延迟付款。非由政府完全信用支持的政府支持实体或联邦机构证券，不属于美国政府直接发行或担保。


\textbf{抵押及资产支持证券。}基础资产可能是住房贷款、消费者贷款或应收款。除一般债务风险外，借款人付款速度变化影响价值。利率下降时提前还款增多，须以较低收益再投资；利率上升时偿还较慢，资金不能及时再投资高收益工具。借款人可能违约，担保房地产可能不足以在止赎后偿债。其他资产支持证券面临类似风险。


\textbf{流动性风险。}部分持仓在市场动荡时难以或无法出售。没有活跃市场、法律或合同转售限制，或做市人减少，都可能降低流动性。价格可能波动或难以确定。出售可能延误且成本增加，基金可能无法在需要时出售，或只能亏损出售。


美国以外投资。政治社会经济和市场变化、地区冲突及战争、汇率变化、国有化、外汇限制、制裁、税收和价格控制，可降低价值。外国市场可能波动更大、流动性更低，适用不同监管、法律、会计、审计及披露要求，估值更困难。外国税收和交易结算延误也可影响结果，新兴市场风险可能更高。


\subsection*{其他投资风险}
\textbf{利率风险。} 利率变化可能影响基金持有证券的价值及流动性。例如，利率上升时，债务证券价值可能下降；利率下降时，其价值可能上升。较长期限的债务证券通常对利率变化更为敏感，价格波动也可能大于较短期限证券。基金可投资可变利率及浮动利率证券。持有此类证券时，市场利率下降会不利于基金收到的收入及份额净值。虽然此类证券的价值一般比其他债务证券对利率变化不敏感，但如果其利率上调的速度或幅度低于市场利率，其价值仍可能下降。相反，利率下降时，浮动利率证券的价值通常不会上升。短期利率极低时，部分债务证券可能不能维持正收益或正总回报；在利率较低的环境下，利率上升的相关风险较高。


\textbf{国家、地区、行业或部门敞口。} 在基金投资限制允许的范围内，基金可能对某一国家、地区、行业或部门具有较大敞口。相关风险及事件对基金的影响可能因此更大，使其份额净值比没有同等敞口的基金更为波动。例如，对某一国家有较大敞口时，该国不利的社会、经济、监管或其他事件，对基金的影响可能大于对地理分布更广的基金的影响。


\textbf{组合证券出借。} 证券出借涉及风险，包括借出证券不能及时归还或者无法归还，以致影响基金行使代理投票权或完成交易，以及交易对手违约的风险。基金还可能因将出借证券收到的担保品再投资于价值下跌、发生违约或表现不及预期的投资而亏损。


\textbf{网络安全事件。} 网络安全遭到破坏，可能使基金承担运营及信息安全风险。此类事件可能源于蓄意攻击或非故意事件，包括勒索软件攻击、注入计算机病毒或恶意代码、利用程序漏洞未经授权访问数字信息系统、网络或设备，以及针对投资顾问或关联机构网站的拒绝服务等外部攻击，使基金网络服务不能向预期终端用户提供。事件可能造成机密资料被未经授权披露、基金资产或敏感信息被滥用、运营能力中断、持有人无法交易，以及实体基础设施、设备或操作系统损毁。由此可能违反隐私及其他法律，损害声誉，产生补救成本或财务损失。投资顾问、过户代理人、托管人、行政服务商及其他金融中介等第三方服务提供者发生类似事件时，基金亦可能面临相关风险。所投资证券的发行人也可能发生网络安全事件，导致基金投资价值下跌。


\textbf{大额持有人交易风险。} 相对于基金规模，持有人，包括投资顾问管理的其他基金或账户，单独或合计大额购买或赎回份额，可能对基金产生不利影响。例如，投资顾问调整其管理的其他基金及账户的配置，可能引致相对于本基金规模较大的交易。大额赎回可能迫使基金在原本不会出售证券的时点出售组合证券，不利于净值及流动性；大额购买则可能在新资金尚未完成投资期间，使基金维持高于通常水平的现金头寸，不利于业绩。因出售投资实现收益，此类交易还可能提前产生须向持有人分配的应税收入，并增加交易成本。大额赎回还可能使现有费用分摊于较小的资产规模，提高费用率。基金规模较小时，上述风险较高。


除上述主要投资策略外，补充资料还载有其他投资方式，以及主要策略和其他投资方式涉及的其他风险。基金投资结果取决于投资顾问处理上述风险及补充资料所述风险的能力。


\subsection*{基金比较指数}
标普500指数按市值加权，以约500只广泛持有的普通股的投资结果为基础。该指数未经主动管理，其结果包括股息或分配的再投资，不反映销售费、佣金、账户费用、运营支出或美国联邦所得税。


60\%标普500指数／40\%彭博美国综合指数，分别按60\%及40\%的权重组合两指数的累计总回报，并假设每月再平衡。该组合指数未经主动管理，其结果包括股息或分配的再投资，不反映销售费、佣金、账户费用、运营支出或美国联邦所得税。


彭博美国综合指数代表美国投资级固定利率债券市场。该指数未经主动管理，其结果包括股息或分配的再投资，不反映销售费、佣金、账户费用、运营支出或美国联邦所得税。


\subsection*{基金业绩计算}
本招募说明书所列基金业绩均反映股息及资本利得分配的再投资。除另有说明外，基金业绩反映列示期间内有效的费用减免及费用补偿。


\subsection*{组合持仓}
基金组合持仓资料可向基金管理机构查询。有关持仓信息披露的政策及程序，载于补充资料。


\section*{管理与组织}
\subsection*{投资顾问}
管理公司M担任本基金及同系列其他基金的投资顾问，负责管理基金投资组合及业务事务。最近财务年度支付给投资顾问的管理费，占平均净资产的百分比，列于“基金费用及支出”的年度基金运营费用表；详情参见补充资料。基金董事会批准投资顾问及服务协议的依据，载于最近财务期间的年度或半年度N-CSR报告。


投资顾问通过三个股票投资部门及一个固定收益投资部门管理资产。各股票部门独立作出投资决定。


股票投资部门日后可能改组为投资顾问的全资子公司；届时投资顾问仍担任基金投资顾问，日常股票投资管理继续通过一个或多个子公司进行。虽然目前无此计划，固定收益投资部门日后亦可能改组，并受聘提供日常固定收益投资管理服务。投资顾问及其管理的基金已取得美国证券交易委员会的豁免命令，允许在基金董事会批准后，使用管理子公司及关联机构提供日常投资管理服务，包括变更提供服务的管理子公司及关联机构。持有人已批准此安排，但不能保证投资顾问将改组其投资部门或行使豁免命令授予的权限。


\subsection*{多组合经理制度}
投资顾问将基金组合划分为多个部分，由不同经理分别管理。投资分析师亦可对组合的一部分作出投资决定。有关决定须符合基金目标、政策及限制，并接受投资顾问及其投资部门有关投资委员会的监督。


固定收益投资部门的部分资深成员组成组合策略小组，以研究为基础，结合分析师、组合经理及经济学家的意见，就久期、收益率曲线及行业配置等宏观因素确定投资主题。本基金固定收益组合经理作出投资决定时，会参考该小组讨论的投资主题。


下表列示本基金主要组合经理的投资经验及管理职责。


\Needspace{5\baselineskip}
\begingroup\small\setlength{\tabcolsep}{4pt}\renewcommand{\arraystretch}{1.35}
\begin{longtable}{>{\centering\arraybackslash}m{0.12192\dimexpr\linewidth-8\tabcolsep\relax}>{\centering\arraybackslash}m{0.27807\dimexpr\linewidth-8\tabcolsep\relax}>{\centering\arraybackslash}m{0.36418\dimexpr\linewidth-8\tabcolsep\relax}>{\centering\arraybackslash}m{0.23583\dimexpr\linewidth-8\tabcolsep\relax}}
\toprule
\rowcolor{TableHead} \centering\arraybackslash \textbf{组\mbox{合经理}} & \centering\arraybackslash \textbf{投资经验及加入顾问机\mbox{构年份}} & \centering\arraybackslash \textbf{管理本基金起\mbox{始年份}} & \centering\arraybackslash \textbf{管\mbox{理职责}} \\
\midrule\endfirsthead
\toprule
\rowcolor{TableHead} \centering\arraybackslash \textbf{组\mbox{合经理}} & \centering\arraybackslash \textbf{投资经验及加入顾问机\mbox{构年份}} & \centering\arraybackslash \textbf{管理本基金起\mbox{始年份}} & \centering\arraybackslash \textbf{管\mbox{理职责}} \\
\midrule\endhead
\midrule\endfoot
\bottomrule\endlastfoot
\centering\arraybackslash 经理01 & \centering\arraybackslash 1983年从业；1995\mbox{年加入} & \centering\arraybackslash 1999年；1997年起曾任本基金\mbox{分析师} & \centering\arraybackslash 股票及固定收益组\mbox{合管理} \\ \addlinespace[3pt]
\centering\arraybackslash 经理02 & \centering\arraybackslash 2004年从业；2016\mbox{年加入} & \centering\arraybackslash 2018年 & \centering\arraybackslash 固定收益组\mbox{合管理} \\ \addlinespace[3pt]
\centering\arraybackslash 经理03 & \centering\arraybackslash 2005年从业\mbox{及加入} & \centering\arraybackslash 2014年；2006年起曾任本基金\mbox{分析师} & \centering\arraybackslash 股票组\mbox{合管理} \\ \addlinespace[3pt]
\centering\arraybackslash 经理04 & \centering\arraybackslash 1986年从业；1991\mbox{年加入} & \centering\arraybackslash 2005年；2000年起曾任本基金\mbox{分析师} & \centering\arraybackslash 股票组\mbox{合管理} \\ \addlinespace[3pt]
\centering\arraybackslash 经理05 & \centering\arraybackslash 1997年从业；2004\mbox{年加入} & \centering\arraybackslash 2021年 & \centering\arraybackslash 股票组\mbox{合管理} \\ \addlinespace[3pt]
\centering\arraybackslash 经理06 & \centering\arraybackslash 1996年从业；1997\mbox{年加入} & \centering\arraybackslash 2018年 & \centering\arraybackslash 股票组\mbox{合管理} \\ \addlinespace[3pt]
\centering\arraybackslash 经理07 & \centering\arraybackslash 2004年从业；2022\mbox{年加入} & \centering\arraybackslash 2022年 & \centering\arraybackslash 固定收益组\mbox{合管理} \\ \addlinespace[3pt]
\centering\arraybackslash 经理08 & \centering\arraybackslash 1989年从业；2002\mbox{年加入} & \centering\arraybackslash 2016年 & \centering\arraybackslash 固定收益组\mbox{合管理} \\ \addlinespace[3pt]
\centering\arraybackslash 经理09 & \centering\arraybackslash 2001年从业；2005\mbox{年加入} & \centering\arraybackslash 2020年 & \centering\arraybackslash 股票组\mbox{合管理} \\ \addlinespace[3pt]
\centering\arraybackslash 经理10 & \centering\arraybackslash 2000年从业；2011\mbox{年加入} & \centering\arraybackslash 2020年 & \centering\arraybackslash 固定收益组\mbox{合管理} \\ \addlinespace[3pt]
\centering\arraybackslash 经理11 & \centering\arraybackslash 1991年从业\mbox{及加入} & \centering\arraybackslash 2016年 & \centering\arraybackslash 股票组\mbox{合管理} \\
\end{longtable}\endgroup
组合经理报酬、持有本基金证券的情况及其管理的其他账户，载于补充资料。


本招募说明书及补充资料所列的部分权利或服务，可能因您的投资经销商或退休计划记录机构而不能向您提供。详情请向财务专业人员或退休计划记录机构查询。


\section*{持有人信息}
过户代理人提供可随需要变化而调整投资安排的服务。服务可在提前60日书面通知后终止或修改。更详细政策和服务见补充资料及新持有人手册。


\subsection*{购买、转换及出售份额}
过户代理人依法须取得客户或代理人的个人信息以核实身份。资料不足可能不能开户；身份无法核实或发现潜在犯罪活动时，基金和分销机构保留关闭账户或采取合理或法定行动的权利。


购买时应指定基金。未指定时，资金可不投资等待指令，最多三个工作日；代理人不对等待期间收入或升值损失负责。接受投资指令后，按下次确定的净值及适用销售费投资。未取得指令时，按规则在第三工作日投入政府货币市场基金。10,000美元或以下且过去16个月有现金投资记录等特定情形，可能按上次基金、比例及类别处理。只有一个已开基金时，指令完整的款项可在收到当天投入该基金。雇主安排可能适用不同流程。


\subsection*{份额估值}
各类别净值为其占有的投资、现金及其他资产价值，减去分配给该类别的负债，再除以已发行份额数。已实现收入和收益纳入净值至除息日，届时已宣布分红成为基金负债。


通常在纽约证券交易所每个营业日常规收盘时计算一次，通常纽约时间下午4时。交易所提前收盘时，约在该时点估值；若相关市场继续交易，第三方报价按较晚时点提供，可使用这些报价。


股票主要按市场报价，债券主要按第三方定价服务估值。报价不可得或不可靠时采用公允价值程序。例如外国市场收盘后发生重大事件，或债券违约且缺乏市场时，可能按公允价值估值。此程序旨在提供更适当净值并减少短期套利。


基金可能持有外国证券，在基金不估值或不允许买卖的日子，证券价值仍可变化。指令及法律文件齐备后，交易按下次净值办理。收盘后收到的完整订单按下一营业日净值办理。


\subsection*{收费渠道及退休计划类别购买}
收费渠道类别通常通过与分销机构有专门协议的收费投资项目、获准自助经纪平台、注册投资顾问或其他获准中介购买。中介通常另收持续服务费，费用不由基金支付。某些代理执行经纪平台还可另收佣金或其他报酬。


获基金顾问高管批准的机构投资者，包括慈善组织、政府机构、公司和金融中介，也可持有相应机构类别。直接由过户代理人持有和服务的该类账户通常须至少100万美元；收费项目持有的某些类别，其购买最低额目前可豁免。


退休计划类别通常仅供税法第401(a)、403(b)、457条计划、非合资格递延报酬安排及某些员工福利或退休后福利计划。通常须在基金记录上以计划级或综合账户持有。部分退休类别仅通过收费渠道或退休计划中介提供。并非每个经销商都提供全部类别。


除说明书明确例外，退休计划类别通常不提供给普通非退休个人账户、传统或Roth IRA、经纪账户持有的SEP、SARSEPs、SIMPLE IRA及529计划。某些雇主退休安排通过退休记录机构服务并使用收费项目时可适用例外。购买具体由计划管理人或记录机构办理。


通常账户初始最低250美元，追加最低50美元；工资扣款退休账户初始或追加最低25美元。部分最低额可以豁免。低于最低初始额的账户，基金保留按当时净值赎回的权利，具体见补充资料。


\subsection*{份额转换}
除不享有转换特权的类别及其他明确例外，可转换为其他同系列基金的同类份额而不另收销售费。曾无销售费买入的货币市场基金转入其他基金，可能产生适用销售费。类别及账户之间移动投资的限制见补充资料。


转换一般具有与出售及买入相同的税务结果。在应税账户，价值高于成本的份额转换可能产生税负。税收优惠退休账户内转换通常不产生联邦或州所得税资本损益。


\subsection*{如何出售份额}
雇主退休计划。通过计划管理人或记录机构出售。


通过经销商或财务顾问。以经销商名义持有的份额须通过该经销商出售；收费渠道份额须经相应中介办理，可能产生中介收费。


书面申请。须由登记持有人签署。金额超过250,000美元、支付给非登记持有人、发送至非登记地址或过去10日变更的登记地址时，须签名担保。代理人可对任何赎回要求担保；公司、合伙或受托账户可要求额外文件。


电话或传真。每名持有人每天限250,000美元。支票须付给登记持有人，并寄至使用至少10日的登记地址。网上或语音自助赎回每天限125,000美元，同样适用收款人及地址要求。


通常预计接受完整订单后一个工作日发送赎回款，但可能超过一日，依法一般可至七日；有限例外可更久。近期买入后赎回，可等待足够时间确认支票或付款已清算，通常自购买日起七个工作日。


正常情况下从现金等高流动性资产支付；必要时可出售证券、使用托管透支或信用额度、向其他关联管理基金借款，或用组合证券等实物支付。基金信托声明允许董事会决定部分或全部实物支付，同一天不同持有人可取得不同证券或现金。取得实物后出售可能支付经纪成本，并在出售前承担市场及流动性风险。


\subsection*{电话、传真或网上交易}
通常自动取得电子赎回和转换服务，可书面选择不使用，也可恢复。代理人采取合理程序确认拥有账户信息者的指令真实时，客户同意基金、代理人及相关人员不承担使用该服务的损失及费用。未采取合理程序时，代理人或基金可对未经授权或欺诈指令的损失负责。


\subsection*{频繁交易}
基金和分销机构可因任何原因拒绝购买订单。基金不为频繁交易设计；频繁交易可提高成本、降低管理效率并稀释长期持有人价值。实际或潜在损害基金的买入，包括转换，可以拒绝。


过户代理人监测大额交易、相近时间交易、共同所有权或影响下的多个账户交易等，判断是否构成频繁交易。监测程序可为改善识别、退休计划监测和法律合规而调整。


某些活动通常不作为频繁交易，包括可识别的退休计划缴款、贷款和分配，基金之间为满足申赎的交易，社区基金会账户，实物转入、滚存、Roth IRA转换及重定性，以及获准对冲或实物赎回安排。可识别的系统性申赎通常须预先安排特定日期。


代理人可接受中介自己的监测程序，但须认为程序合理执行基金政策。发现潜在频繁交易时可要求中介提供持有人信息或限制交易；中介未适当处理时可终止交易资格。不保证防止所有频繁交易；所有交易仍受限制滥用交易的一般权利约束。


\section*{分配及税务}
基金通常在3、6、9、12月分红；资本利得（如有）通常在12月和6月分配。分配时每份净值减少相应金额。持有人可选择现金或再投资本基金及其他同系列基金；退休计划及教育储蓄类别自动再投资。


联邦税务上，股息和短期资本利得通常作为普通收入，符合持有期等条件的合资格股息可能适用较低税率，净长期资本利得按长期资本利得处理。无论现金或再投资，通常在分配时应税。税收优惠退休或教育账户中自动再投资通常不在当时产生联邦或州所得税。


赎回及转换可产生资本损益，为出售所得与含销售费的成本之差。税收优惠退休账户内部转换一般不产生联邦或州资本损益；有限例外外，退休账户提款通常按普通收入纳税。具体请咨询税务顾问。


\section*{选择份额类别}
每类投资同一组合，但销售费及费用结构不同。收费渠道中的不同类别可分别有12b-1和子过户服务费、仅有子过户费，或不收这些费用，以不同方式支付顾问平台成本。退休类别有不同12b-1和记录管理费，便于计划选择与投资及参与者记录服务成本相适应的类别。


选择时应考虑预期持有期限、投资额、各类总费用、销售费减免资格、转换灵活性及近期分配需要。可用服务和类别会随购买渠道不同。某个类别在基金层存在，并不表示某个渠道必然提供。


\section*{分销计划}
本基金按12b-1计划，为适用份额类别支付主要用于销售份额的活动费用；有关费用类别须由董事会预先批准。P类的计划付款上限为平均每日净资产的年化1.00\%，其中最高0.25\%可用于向合资格经销商支付提供持有人服务的服务费，余款可用于分销支出。


最近财务年度实际支付的12b-1费率列于“基金费用及支出”的年度基金运营费用表。由于这些费用持续从基金资产支付，长期累计可能高于其他销售费或服务费，并降低投资回报。


\section*{向经销商支付的其他报酬}
分销机构以自身费用向投资经销商支付额外报酬。按资产规模计算的合资格付款，一般采用一致公式，最高不超过归属于该经销商的合资格同系列基金资产的0.035\%。合资格资产不包括退休计划类份额、Q类份额、IRA中持有的收费渠道类份额及某些退休账户份额。经销商可指示分销机构排除其他资产。


分销机构亦就渠道平台、清算、托管和记录管理、账户维护、报表制作、交易处理、运营改进、培训、教育、市场推广、交易费用支持、技术及数据等服务，向某些金融中介及其他机构付款。


分销机构每年决定是否继续支付额外报酬。收款机构须确认付款用途，通常须完成平台尽职调查，不向财务专业人员、分支经理或相关人员提供偏向某一获准基金集团的财务激励，为客户提供个别化建议的机会，向分销机构提供专业人员和产品平台的适当接触机会，并与过户代理人配合提高运营效率及联系持有人。


投资顾问、分销机构或其他关联方，如向经销商支付不同金额的额外报酬或激励，可能使经销商及其专业人员具有推荐某一共同基金而非其他共同基金或投资的财务动机，构成潜在利益冲突。您应咨询财务专业人员，并仔细阅读其机构关于所收报酬的披露。


\section*{基金支出}
市场波动期间，基金资产可能大幅下降，导致年度运营费用占投资价值的百分比高于“基金费用及支出”所列数值。


所有类别的“其他费用”包括投资顾问及关联机构提供行政服务的费用。有关服务协助为持有人提供非分销服务的第三方，包括提供基金及影响基金投资的市场发展的深入资料，以及协调、监控及监督第三方持有人服务。行政服务协议允许对全部类别收取最高0.05\%的行政服务费。投资顾问目前按平均每日净资产的年化0.03\%收取该费用，可按前述约定提高。


“其他费用”亦包括托管、法律、过户代理，以及适用时的子过户代理和记录管理付款，及各类别适用的其他费用。


\subsection*{子过户代理及记录管理费用}
第三方，包括投资顾问的关联机构，可代替过户代理人为某些持有人账户提供子过户代理、记录管理或持有人服务，并收取相应付款。金额因份额类别及服务而异，通常为每账户3—18美元。Q类不承担子过户代理或记录管理费用。


雇主退休计划的子过户代理及记录管理付款，因计划选择的类别而异。P类向提供此类服务的机构支付的最高费用为资产的0.35\%。


\section*{财务摘要}
财务摘要旨在帮助您了解基金过去五个财务年度的投资结果。部分资料反映有关类别单一份额的财务结果。总回报假设所有股息及资本利得分配均再投资。下列财务资料已经独立注册会计师审计。审计报告及基金财务报表载于补充资料，可应要求取得。


\subsection*{P类份额}
每份财务数据以美元表示；各年度截至12月31日。


\Needspace{5\baselineskip}
\begingroup\small\setlength{\tabcolsep}{4pt}\renewcommand{\arraystretch}{1.35}
\begin{longtable}{>{\centering\arraybackslash}m{0.16391\dimexpr\linewidth-12\tabcolsep\relax}>{\centering\arraybackslash}m{0.16391\dimexpr\linewidth-12\tabcolsep\relax}>{\centering\arraybackslash}m{0.16391\dimexpr\linewidth-12\tabcolsep\relax}>{\centering\arraybackslash}m{0.18045\dimexpr\linewidth-12\tabcolsep\relax}>{\centering\arraybackslash}m{0.16391\dimexpr\linewidth-12\tabcolsep\relax}>{\centering\arraybackslash}m{0.16391\dimexpr\linewidth-12\tabcolsep\relax}}
\toprule
\rowcolor{TableHead} \centering\arraybackslash \textbf{年度} & \centering\arraybackslash \textbf{年\mbox{初净值}} & \centering\arraybackslash \textbf{净投\mbox{资收入}} & \centering\arraybackslash \textbf{证券已实现及未实现\mbox{净损益}} & \centering\arraybackslash \textbf{投资运\mbox{作合计}} & \centering\arraybackslash \textbf{年\mbox{末净值}} \\
\midrule\endfirsthead
\toprule
\rowcolor{TableHead} \centering\arraybackslash \textbf{年度} & \centering\arraybackslash \textbf{年\mbox{初净值}} & \centering\arraybackslash \textbf{净投\mbox{资收入}} & \centering\arraybackslash \textbf{证券已实现及未实现\mbox{净损益}} & \centering\arraybackslash \textbf{投资运\mbox{作合计}} & \centering\arraybackslash \textbf{年\mbox{末净值}} \\
\midrule\endhead
\midrule\endfoot
\bottomrule\endlastfoot
\centering\arraybackslash 2025 & \centering\arraybackslash 34.05 & \centering\arraybackslash .48 & \centering\arraybackslash 5.44 & \centering\arraybackslash 5.92 & \centering\arraybackslash 37.15 \\ \addlinespace[3pt]
\centering\arraybackslash 2024 & \centering\arraybackslash 31.74 & \centering\arraybackslash .50 & \centering\arraybackslash 4.02 & \centering\arraybackslash 4.52 & \centering\arraybackslash 34.05 \\ \addlinespace[3pt]
\centering\arraybackslash 2023 & \centering\arraybackslash 28.55 & \centering\arraybackslash .48 & \centering\arraybackslash 3.24 & \centering\arraybackslash 3.72 & \centering\arraybackslash 31.74 \\ \addlinespace[3pt]
\centering\arraybackslash 2022 & \centering\arraybackslash 33.22 & \centering\arraybackslash .40 & \centering\arraybackslash (4.64) & \centering\arraybackslash (4.24) & \centering\arraybackslash 28.55 \\ \addlinespace[3pt]
\centering\arraybackslash 2021 & \centering\arraybackslash 29.99 & \centering\arraybackslash .23 & \centering\arraybackslash 4.19 & \centering\arraybackslash 4.42 & \centering\arraybackslash 33.22 \\
\end{longtable}\endgroup
\Needspace{5\baselineskip}
\begingroup\small\setlength{\tabcolsep}{4pt}\renewcommand{\arraystretch}{1.35}
\begin{longtable}{>{\centering\arraybackslash}m{0.25\dimexpr\linewidth-8\tabcolsep\relax}>{\centering\arraybackslash}m{0.25\dimexpr\linewidth-8\tabcolsep\relax}>{\centering\arraybackslash}m{0.25\dimexpr\linewidth-8\tabcolsep\relax}>{\centering\arraybackslash}m{0.25\dimexpr\linewidth-8\tabcolsep\relax}}
\toprule
\rowcolor{TableHead} \centering\arraybackslash \textbf{年度} & \centering\arraybackslash \textbf{净投资收\mbox{入股息}} & \centering\arraybackslash \textbf{资本利\mbox{得分配}} & \centering\arraybackslash \textbf{股息及分\mbox{配合计}} \\
\midrule\endfirsthead
\toprule
\rowcolor{TableHead} \centering\arraybackslash \textbf{年度} & \centering\arraybackslash \textbf{净投资收\mbox{入股息}} & \centering\arraybackslash \textbf{资本利\mbox{得分配}} & \centering\arraybackslash \textbf{股息及分\mbox{配合计}} \\
\midrule\endhead
\midrule\endfoot
\bottomrule\endlastfoot
\centering\arraybackslash 2025 & \centering\arraybackslash (.50) & \centering\arraybackslash (2.32) & \centering\arraybackslash (2.82) \\ \addlinespace[3pt]
\centering\arraybackslash 2024 & \centering\arraybackslash (.46) & \centering\arraybackslash (1.75) & \centering\arraybackslash (2.21) \\ \addlinespace[3pt]
\centering\arraybackslash 2023 & \centering\arraybackslash (.53) & \centering\arraybackslash — & \centering\arraybackslash (.53) \\ \addlinespace[3pt]
\centering\arraybackslash 2022 & \centering\arraybackslash (.25) & \centering\arraybackslash (.18) & \centering\arraybackslash (.43) \\ \addlinespace[3pt]
\centering\arraybackslash 2021 & \centering\arraybackslash (.15) & \centering\arraybackslash (1.04) & \centering\arraybackslash (1.19) \\
\end{longtable}\endgroup
\Needspace{5\baselineskip}
\begingroup\small\setlength{\tabcolsep}{4pt}\renewcommand{\arraystretch}{1.35}
\begin{longtable}{>{\centering\arraybackslash}m{0.19601\dimexpr\linewidth-10\tabcolsep\relax}>{\centering\arraybackslash}m{0.19601\dimexpr\linewidth-10\tabcolsep\relax}>{\centering\arraybackslash}m{0.20269\dimexpr\linewidth-10\tabcolsep\relax}>{\centering\arraybackslash}m{0.19601\dimexpr\linewidth-10\tabcolsep\relax}>{\centering\arraybackslash}m{0.20928\dimexpr\linewidth-10\tabcolsep\relax}}
\toprule
\rowcolor{TableHead} \centering\arraybackslash \textbf{年度} & \centering\arraybackslash \textbf{总回报（\%）} & \centering\arraybackslash \textbf{年末净资产（百\mbox{万美元）}} & \centering\arraybackslash \textbf{费用／平均净资产（\%）} & \centering\arraybackslash \textbf{净收入／平均净资产（\%）} \\
\midrule\endfirsthead
\toprule
\rowcolor{TableHead} \centering\arraybackslash \textbf{年度} & \centering\arraybackslash \textbf{总回报（\%）} & \centering\arraybackslash \textbf{年末净资产（百\mbox{万美元）}} & \centering\arraybackslash \textbf{费用／平均净资产（\%）} & \centering\arraybackslash \textbf{净收入／平均净资产（\%）} \\
\midrule\endhead
\midrule\endfoot
\bottomrule\endlastfoot
\centering\arraybackslash 2025 & \centering\arraybackslash 17.58 & \centering\arraybackslash 1,160 & \centering\arraybackslash 1.33 & \centering\arraybackslash 1.34 \\ \addlinespace[3pt]
\centering\arraybackslash 2024 & \centering\arraybackslash 14.08 & \centering\arraybackslash 1,147 & \centering\arraybackslash 1.32 & \centering\arraybackslash 1.47 \\ \addlinespace[3pt]
\centering\arraybackslash 2023 & \centering\arraybackslash 13.11 & \centering\arraybackslash 1,106 & \centering\arraybackslash 1.33 & \centering\arraybackslash 1.63 \\ \addlinespace[3pt]
\centering\arraybackslash 2022 & \centering\arraybackslash (12.78) & \centering\arraybackslash 1,070 & \centering\arraybackslash 1.34 & \centering\arraybackslash 1.34 \\ \addlinespace[3pt]
\centering\arraybackslash 2021 & \centering\arraybackslash 14.86 & \centering\arraybackslash 1,327 & \centering\arraybackslash 1.33 & \centering\arraybackslash .73 \\
\end{longtable}\endgroup
\subsection*{Q类份额}
每份财务数据以美元表示；各年度截至12月31日。


\Needspace{5\baselineskip}
\begingroup\small\setlength{\tabcolsep}{4pt}\renewcommand{\arraystretch}{1.35}
\begin{longtable}{>{\centering\arraybackslash}m{0.16391\dimexpr\linewidth-12\tabcolsep\relax}>{\centering\arraybackslash}m{0.16391\dimexpr\linewidth-12\tabcolsep\relax}>{\centering\arraybackslash}m{0.16391\dimexpr\linewidth-12\tabcolsep\relax}>{\centering\arraybackslash}m{0.18045\dimexpr\linewidth-12\tabcolsep\relax}>{\centering\arraybackslash}m{0.16391\dimexpr\linewidth-12\tabcolsep\relax}>{\centering\arraybackslash}m{0.16391\dimexpr\linewidth-12\tabcolsep\relax}}
\toprule
\rowcolor{TableHead} \centering\arraybackslash \textbf{年度} & \centering\arraybackslash \textbf{年\mbox{初净值}} & \centering\arraybackslash \textbf{净投\mbox{资收入}} & \centering\arraybackslash \textbf{证券已实现及未实现\mbox{净损益}} & \centering\arraybackslash \textbf{投资运\mbox{作合计}} & \centering\arraybackslash \textbf{年\mbox{末净值}} \\
\midrule\endfirsthead
\toprule
\rowcolor{TableHead} \centering\arraybackslash \textbf{年度} & \centering\arraybackslash \textbf{年\mbox{初净值}} & \centering\arraybackslash \textbf{净投\mbox{资收入}} & \centering\arraybackslash \textbf{证券已实现及未实现\mbox{净损益}} & \centering\arraybackslash \textbf{投资运\mbox{作合计}} & \centering\arraybackslash \textbf{年\mbox{末净值}} \\
\midrule\endhead
\midrule\endfoot
\bottomrule\endlastfoot
\centering\arraybackslash 2025 & \centering\arraybackslash 34.32 & \centering\arraybackslash .88 & \centering\arraybackslash 5.48 & \centering\arraybackslash 6.36 & \centering\arraybackslash 37.47 \\ \addlinespace[3pt]
\centering\arraybackslash 2024 & \centering\arraybackslash 31.97 & \centering\arraybackslash .88 & \centering\arraybackslash 4.05 & \centering\arraybackslash 4.93 & \centering\arraybackslash 34.32 \\ \addlinespace[3pt]
\centering\arraybackslash 2023 & \centering\arraybackslash 28.75 & \centering\arraybackslash .81 & \centering\arraybackslash 3.26 & \centering\arraybackslash 4.07 & \centering\arraybackslash 31.97 \\ \addlinespace[3pt]
\centering\arraybackslash 2022 & \centering\arraybackslash 33.45 & \centering\arraybackslash .73 & \centering\arraybackslash (4.67) & \centering\arraybackslash (3.94) & \centering\arraybackslash 28.75 \\ \addlinespace[3pt]
\centering\arraybackslash 2021 & \centering\arraybackslash 30.18 & \centering\arraybackslash .59 & \centering\arraybackslash 4.22 & \centering\arraybackslash 4.81 & \centering\arraybackslash 33.45 \\
\end{longtable}\endgroup
\Needspace{5\baselineskip}
\begingroup\small\setlength{\tabcolsep}{4pt}\renewcommand{\arraystretch}{1.35}
\begin{longtable}{>{\centering\arraybackslash}m{0.25\dimexpr\linewidth-8\tabcolsep\relax}>{\centering\arraybackslash}m{0.25\dimexpr\linewidth-8\tabcolsep\relax}>{\centering\arraybackslash}m{0.25\dimexpr\linewidth-8\tabcolsep\relax}>{\centering\arraybackslash}m{0.25\dimexpr\linewidth-8\tabcolsep\relax}}
\toprule
\rowcolor{TableHead} \centering\arraybackslash \textbf{年度} & \centering\arraybackslash \textbf{净投资收\mbox{入股息}} & \centering\arraybackslash \textbf{资本利\mbox{得分配}} & \centering\arraybackslash \textbf{股息及分\mbox{配合计}} \\
\midrule\endfirsthead
\toprule
\rowcolor{TableHead} \centering\arraybackslash \textbf{年度} & \centering\arraybackslash \textbf{净投资收\mbox{入股息}} & \centering\arraybackslash \textbf{资本利\mbox{得分配}} & \centering\arraybackslash \textbf{股息及分\mbox{配合计}} \\
\midrule\endhead
\midrule\endfoot
\bottomrule\endlastfoot
\centering\arraybackslash 2025 & \centering\arraybackslash (.89) & \centering\arraybackslash (2.32) & \centering\arraybackslash (3.21) \\ \addlinespace[3pt]
\centering\arraybackslash 2024 & \centering\arraybackslash (.83) & \centering\arraybackslash (1.75) & \centering\arraybackslash (2.58) \\ \addlinespace[3pt]
\centering\arraybackslash 2023 & \centering\arraybackslash (.85) & \centering\arraybackslash — & \centering\arraybackslash (.85) \\ \addlinespace[3pt]
\centering\arraybackslash 2022 & \centering\arraybackslash (.58) & \centering\arraybackslash (.18) & \centering\arraybackslash (.76) \\ \addlinespace[3pt]
\centering\arraybackslash 2021 & \centering\arraybackslash (.50) & \centering\arraybackslash (1.04) & \centering\arraybackslash (1.54) \\
\end{longtable}\endgroup
\Needspace{5\baselineskip}
\begingroup\small\setlength{\tabcolsep}{4pt}\renewcommand{\arraystretch}{1.35}
\begin{longtable}{>{\centering\arraybackslash}m{0.19601\dimexpr\linewidth-10\tabcolsep\relax}>{\centering\arraybackslash}m{0.19601\dimexpr\linewidth-10\tabcolsep\relax}>{\centering\arraybackslash}m{0.20269\dimexpr\linewidth-10\tabcolsep\relax}>{\centering\arraybackslash}m{0.19601\dimexpr\linewidth-10\tabcolsep\relax}>{\centering\arraybackslash}m{0.20928\dimexpr\linewidth-10\tabcolsep\relax}}
\toprule
\rowcolor{TableHead} \centering\arraybackslash \textbf{年度} & \centering\arraybackslash \textbf{总回报（\%）} & \centering\arraybackslash \textbf{年末净资产（百\mbox{万美元）}} & \centering\arraybackslash \textbf{费用／平均净资产（\%）} & \centering\arraybackslash \textbf{净收入／平均净资产（\%）} \\
\midrule\endfirsthead
\toprule
\rowcolor{TableHead} \centering\arraybackslash \textbf{年度} & \centering\arraybackslash \textbf{总回报（\%）} & \centering\arraybackslash \textbf{年末净资产（百\mbox{万美元）}} & \centering\arraybackslash \textbf{费用／平均净资产（\%）} & \centering\arraybackslash \textbf{净收入／平均净资产（\%）} \\
\midrule\endhead
\midrule\endfoot
\bottomrule\endlastfoot
\centering\arraybackslash 2025 & \centering\arraybackslash 18.84 & \centering\arraybackslash 15,026 & \centering\arraybackslash .25 & \centering\arraybackslash 2.41 \\ \addlinespace[3pt]
\centering\arraybackslash 2024 & \centering\arraybackslash 15.32 & \centering\arraybackslash 12,585 & \centering\arraybackslash .25 & \centering\arraybackslash 2.55 \\ \addlinespace[3pt]
\centering\arraybackslash 2023 & \centering\arraybackslash 14.34 & \centering\arraybackslash 10,713 & \centering\arraybackslash .25 & \centering\arraybackslash 2.71 \\ \addlinespace[3pt]
\centering\arraybackslash 2022 & \centering\arraybackslash (11.81) & \centering\arraybackslash 9,501 & \centering\arraybackslash .25 & \centering\arraybackslash 2.45 \\ \addlinespace[3pt]
\centering\arraybackslash 2021 & \centering\arraybackslash 16.13 & \centering\arraybackslash 10,596 & \centering\arraybackslash .25 & \centering\arraybackslash 1.82 \\
\end{longtable}\endgroup
\subsection*{组合周转率}
\Needspace{5\baselineskip}
\begingroup\small\setlength{\tabcolsep}{4pt}\renewcommand{\arraystretch}{1.35}
\begin{longtable}{>{\centering\arraybackslash}m{0.33969\dimexpr\linewidth-12\tabcolsep\relax}>{\centering\arraybackslash}m{0.13206\dimexpr\linewidth-12\tabcolsep\relax}>{\centering\arraybackslash}m{0.13206\dimexpr\linewidth-12\tabcolsep\relax}>{\centering\arraybackslash}m{0.13206\dimexpr\linewidth-12\tabcolsep\relax}>{\centering\arraybackslash}m{0.13206\dimexpr\linewidth-12\tabcolsep\relax}>{\centering\arraybackslash}m{0.13206\dimexpr\linewidth-12\tabcolsep\relax}}
\toprule
\rowcolor{TableHead} \centering\arraybackslash \textbf{截至12月31日\mbox{的年度}} & \centering\arraybackslash \textbf{2025} & \centering\arraybackslash \textbf{2024} & \centering\arraybackslash \textbf{2023} & \centering\arraybackslash \textbf{2022} & \centering\arraybackslash \textbf{2021} \\
\midrule\endfirsthead
\toprule
\rowcolor{TableHead} \centering\arraybackslash \textbf{截至12月31日\mbox{的年度}} & \centering\arraybackslash \textbf{2025} & \centering\arraybackslash \textbf{2024} & \centering\arraybackslash \textbf{2023} & \centering\arraybackslash \textbf{2022} & \centering\arraybackslash \textbf{2021} \\
\midrule\endhead
\midrule\endfoot
\bottomrule\endlastfoot
\centering\arraybackslash 不包括抵押贷款美元滚\mbox{动交易} & \centering\arraybackslash 50\% & \centering\arraybackslash 44\% & \centering\arraybackslash 42\% & \centering\arraybackslash 52\% & \centering\arraybackslash 53\% \\ \addlinespace[3pt]
\centering\arraybackslash 包括抵押贷款美元滚\mbox{动交易} & \centering\arraybackslash 106\% & \centering\arraybackslash 131\% & \centering\arraybackslash 171\% & \centering\arraybackslash 157\% & \centering\arraybackslash 158\% \\
\end{longtable}\endgroup
\textbf{财务摘要附注。} 每份投资运作数据按平均已发行份额计算。总回报不包括任何适用销售费用及递延销售费。费用比率不包括中央基金的费用；本基金间接承担其所投资中央基金的相应费用。


组合周转率不包括基金对中央基金的组合交易。2025年数据不包括实物交易（如有）；2021—2024年数据包括因实物赎回出售证券的价值（如有）。如排除2021年实物赎回所涉出售证券，不包括及包括抵押贷款美元滚动交易的组合周转率分别为46\%及152\%。

\end{document}